% !TITLE 短歌行·其一
% !AUTHOR 曹操
% !DYNASTY 魏晋
% !GENRE 四言古诗
% !ORDER 4
% !SOURCE 《文选》卷二十七，乐府
% !EDITION 《先秦汉魏晋南北朝诗》，逯钦立辑校，中华书局
% !EDITION_STATUS pending

\begin{poem}
对酒当歌，人生几何！
譬如朝露，去日苦多\textsuperscript{1}。
慨当以慷\textsuperscript{2}，忧思难忘。
何以解忧？唯有杜康\textsuperscript{3}。
青青子衿\textsuperscript{4}，悠悠我心。
但为君故，沉吟\textsuperscript{5}至今。
呦呦鹿鸣，食野之苹\textsuperscript{6}。
我有嘉宾，鼓瑟吹笙。
明明如月，何时可掇\textsuperscript{7}？
忧从中来，不可断绝。
越陌度阡\textsuperscript{8}，枉用相存\textsuperscript{9}。
契阔谈讌\textsuperscript{10}，心念旧恩。
月明星稀，乌鹊南飞。
绕树三匝\textsuperscript{11}，何枝可依？
山不厌高，海不厌深\textsuperscript{12}。
周公吐哺\textsuperscript{13}，天下归心。
\end{poem}

\begin{annotations}
\notetext{1}{去日苦多：苦于逝去的日子太多。苦，苦于。}
\notetext{2}{慨当以慷：即“慷慨”，意气激昂。当以，语助词，无实义。}
\notetext{3}{杜康：相传为最早酿酒的人，这里代指酒。}
\notetext{4}{青青子衿，悠悠我心：语出《诗经·郑风·子衿》。衿（jīn），衣领。原诗写思念之情，这里借指对贤才的思慕。}
\notetext{5}{沉吟：低声吟咏，指反复思念。}
\notetext{6}{呦呦鹿鸣，食野之苹：语出《诗经·小雅·鹿鸣》。呦呦，鹿鸣声。苹，艾蒿。原诗写宴饮嘉宾，这里借指招纳贤才。}
\notetext{7}{掇：拾取。掇（duō）。}
\notetext{8}{越陌度阡：穿过纵横的田间小路。南北为阡，东西为陌。指贤才远道来投。}
\notetext{9}{枉用相存：屈尊前来问候。枉，屈尊；存，问候、探望。}
\notetext{10}{契阔谈讌：久别重逢，设宴谈心。契阔，久别。讌，同“宴”。}
\notetext{11}{匝：周，圈。匝（zā）。}
\notetext{12}{山不厌高，海不厌深：山不嫌高，海不嫌深。厌，满足。比喻招纳贤才越多越好。}
\notetext{13}{周公吐哺：周公吃饭时多次吐出口中的食物，起身接待来访的士人。《史记·鲁周公世家》载周公自言“一沐三捉发，一饭三吐哺，起以待士，犹恐失天下之贤人”。诗中借以自比，表明求贤之心。}
\end{annotations}

\begin{appreciation}
《短歌行》是曹操四言诗的代表作之一，也是中国诗里写求贤最有名的篇章之一。“短歌”与“长歌”相对，指歌声短促，本是乐府旧题。

全诗三十二句，一路都在问。“人生几何”“何以解忧”“何时可掇”“何枝可依”，四个问句由远及近，把一个坐拥酒宴的人问得坐立不安。

第一问对着自己。“对酒当歌，人生几何！譬如朝露，去日苦多。”满座酒宴，本该放声高歌，脱口而出的却是人生能有多久。汉乐府挽歌《薤露》唱“薤上露，何易晞”，说的正是人命危浅，曹操把这个比喻接了过来。落到他笔下，忧的不是享乐不够，是功业还没有做成。

第二问对着酒。“慨当以慷，忧思难忘。何以解忧？唯有杜康。”胸中激荡着慷慨之气，忧思却怎么也忘不掉，只好去问酒。酒压得住一时的忧，压不住根本的忧。

第三问才揭开这个根本。“青青子衿，悠悠我心”，这是《诗经·郑风·子衿》的成句，原诗写一个人想念另一个人。衿是衣领，青衿是学子的服色。曹操没有明说想谁，只接上一句“但为君故，沉吟至今”，为了你，我低声吟咏到现在。这个“君”不是某一个人，是所有还没有来到他身边的贤才。下面“呦呦鹿鸣，食野之苹。我有嘉宾，鼓瑟吹笙”，又是《诗经》，这回是《小雅·鹿鸣》，写主人鼓瑟吹笙款待嘉宾。前面读过这首宴饮诗，曹操整段搬来，等于把一场现成的宴席摆在贤才面前：你来，我就这样待你。

可宴席还是空的。“明明如月，何时可掇？”贤才像天上的明月，什么时候才能摘到手里。“忧从中来，不可断绝。”上一段才说“我有嘉宾”，这一段就问何时可掇，一近一远，求贤的心在两头之间来回拉扯。于是他想象贤才穿过田野赶来，“越陌度阡，枉用相存。契阔谈讌，心念旧恩”，风尘仆仆，坐下便能抵足而谈。

第四问落到天上。“月明星稀，乌鹊南飞。绕树三匝，何枝可依？”月光明亮，星光稀疏，乌鹊向南飞去，绕树飞了一圈又一圈，哪一枝才可以栖身。这几句是眼前景，也是天下事，乱世里的贤才就像这树上的乌鹊，飞来飞去，找不到可以托身的主君。曹操替他们问出这一句，其实是在替自己招手。

问到这里，不再问了。“山不厌高，海不厌深。周公吐哺，天下归心。”山不嫌自己高，海不嫌自己深，所以能成其高、成其深。周公吃一顿饭，几次把嘴里的食物吐出来，起身接待来访的士人，天下的人心便都归向他。这是表态，也是承诺：我容得下多少人，天下就有多少人肯来。

《诗经》以后，四言诗在文人笔下渐渐沉寂，到曹操手里忽然又活了。钟嵘《诗品》评曹操“古直，甚有悲凉之句”，句子是直的，底色是悲凉的，悲凉底下压着一股不肯认输的劲。

观沧海向外看，龟虽寿向内看，短歌行向人看。三首诗，三种目光，落到同一个人身上。酒宴终有散的时候，天下却还没有归心。
\end{appreciation}
